\section{Предлагаемый метод}

В настоящей главе описан метод обнаружения каузальных зависимостей между парами оптимизирующих проходов LLVM на основе рандомизированного эксперимента. Метод включает: постановку задачи каузального анализа в терминах модели потенциальных исходов, обоснование применимости рандомизированного эксперимента, этапы вычисления среднего эффекта воздействия (ATE) для всех пар проходов с коррекцией множественного тестирования, построение каузального графа на трёх уровнях агрегации, оценку гетерогенности эффектов (CATE), а также сравнение полученного каузального графа с корреляционным графом, построенным на тех же данных.

\subsection{Постановка задачи каузального анализа}

\subsubsection{Единица наблюдения, воздействие и исход}

Единицей наблюдения в работе является \emph{прогон} --- результат применения некоторой случайной последовательности оптимизирующих проходов к заданному бенчмарку. Для каждого прогона~$i$ известны: программа~$P_i$ (один из 517~бенчмарков), последовательность проходов~$\pi_i$ (упорядоченное подмножество из 89~доступных проходов) и итоговое значение метрики --- число IR-инструкций после применения~$\pi_i$.

Для каждой пары проходов $(A, B)$ из $\binom{89}{2} = 3916$~неупорядоченных пар исследуется каузальный вопрос:
\begin{quote}
    \emph{Влияет ли порядок проходов $A$ и $B$ внутри последовательности на итоговое значение метрики качества кода?}
\end{quote}

В терминах модели потенциальных исходов (раздел~3.2) для каждой пары вводятся:

\begin{itemize}
    \item \textbf{Бинарное воздействие} $T \in \{0, 1\}$: переменная, кодирующая порядок проходов $A$ и $B$ в последовательности (точная формулировка зависит от выбранного определения воздействия, см.~ниже).
    \item \textbf{Исход} $Y$: нормированное улучшение метрики, $Y = (M_{\text{raw}} - M_\pi) / M_{\text{raw}}$, где $M_\pi$ --- число IR-инструкций после применения последовательности~$\pi$, а $M_{\text{raw}}$ --- baseline без оптимизаций.
    \item \textbf{Ковариаты}: характеристики прогона, влияющие на исход и потенциально коррелированные с воздействием (присутствие и позиции остальных 87~проходов, длина последовательности, характеристики программы).
\end{itemize}

Каузальный эффект пары $(A, B)$ определяется как $\text{ATE}(A \to B) = \mathbb{E}[Y(1) - Y(0)]$, где $Y(1)$ и $Y(0)$ --- потенциальные исходы при $T = 1$ и $T = 0$ соответственно. В силу рандомизации последовательностей (см. раздел~4.1.3) истинный ATE совпадает с наивной оценкой $\mathbb{E}[Y \mid T = 1] - \mathbb{E}[Y \mid T = 0]$ и интерпретируется как ожидаемое изменение доли устранённых инструкций при изменении порядка~$A$ и~$B$.

\subsubsection{Обоснование выбора нормированного исхода}

Выбор $Y = (M_{\text{raw}} - M_\pi) / M_{\text{raw}}$ обусловлен следующими соображениями:

\begin{itemize}
    \item \textbf{Сопоставимость между бенчмарками.} Абсолютное число устранённых инструкций несравнимо: на крупном бенчмарке (60\,000~инструкций) уменьшение на 100 несущественно, на мелком (140~инструкций) --- значительно. Нормировка на baseline переводит исход в единую безразмерную шкалу для всех 517~бенчмарков.

    \item \textbf{Нормировка на raw, а не на \texttt{-Oz}.} Использование \texttt{-Oz} в качестве знаменателя приводит к зависимости масштаба исхода от качества стандартного пайплайна: если \texttt{-Oz} на конкретном бенчмарке уже близок к оптимуму, любые изменения относительно него выглядят малыми; raw представляет собой чистый baseline без оптимизаций и позволяет единообразно интерпретировать долю улучшения.

    \item \textbf{Прямая интерпретация ATE.} Поскольку $Y$ --- доля улучшения, оценённое значение ATE интерпретируется напрямую. Например, ATE~$= 0{,}039$ означает: <<порядок $A \to B$ улучшает метрику в среднем на 3{,}9~процентных пункта от baseline относительно порядка $B \to A$>>. Это позволяет сравнивать ATE между парами без дополнительной нормировки на стандартное отклонение (см. раздел~3.4.3).
\end{itemize}

\subsubsection{Применимость рандомизированного эксперимента}

В работе случайные последовательности проходов генерируются по фиксированной схеме: для каждого прогона длина $L$~выбирается равномерно из $\{10, 11, \dots, 40\}$, проходы выбираются без повторений из 89~доступных, и расставляются в случайном порядке (равномерное распределение по перестановкам). Параметры выборки одинаковы для всех бенчмарков; на каждый бенчмарк приходится по 10\,000~прогонов.

Для произвольной пары проходов $(A, B)$ относительный порядок $A$ и $B$ в случайной последовательности~$\pi$ не зависит от присутствия и позиций остальных проходов. Это следует из симметрии схемы генерации относительно $A$ и $B$: при условии присутствия обоих проходов оба упорядочения равновероятны ($\mathbb{P}\{A \text{ раньше } B\} = 1/2$), и эта симметрия сохраняется при произвольной фиксации присутствия и позиций остальных элементов~$\pi$.

Следовательно, бинарная переменная воздействия~$T_{AB}^{(1)}$, определяемая через относительный порядок, удовлетворяет условию независимости от потенциальных исходов: $\{Y(1), Y(0)\} \perp T_{AB}^{(1)}$. В соответствии с разделом~3.2.3, наивная оценка $\widehat{\text{ATE}} = \bar{Y}_1 - \bar{Y}_0$ является несмещённой оценкой каузального эффекта.

Это свойство устраняет необходимость в алгоритмах поиска структуры (PC-алгоритм, FCI и др.): ковариаты сбалансированы рандомизацией по построению, и дополнительная корректировка на конфаундеры не требуется.

\subsubsection{Две формулировки воздействия}

Каузальный вопрос о порядке пары проходов допускает несколько естественных формулировок бинарного воздействия. В работе рассматриваются две.

\paragraph{Вариант 1 (относительный порядок).} Для пары $(A, B)$ воздействие определяется только взаимным расположением проходов в последовательности:
\begin{equation*}
    T_{AB}^{(1)} = \begin{cases}
        1, & \text{если } A \text{ встречается раньше } B \text{ в } \pi; \\
        0, & \text{если } B \text{ встречается раньше } A; \\
        \text{не определено}, & \text{если хотя бы один из проходов отсутствует в } \pi.
    \end{cases}
\end{equation*}

Прогоны, в которых отсутствует хотя бы один из проходов пары, исключаются из выборки для данной пары. При параметрах эксперимента ($L \in [10, 40]$ из 89~проходов) доля прогонов с одновременным присутствием обоих проходов составляет около 8~\%, что даёт около 400\,000~наблюдений на пару из 5{,}17~миллиона прогонов.

Вариант~1 улавливает \emph{глобальные} зависимости: эффект сохраняется, даже если между $A$ и $B$ стоит несколько других проходов. Содержательная интерпретация --- <<подготовка IR проходом $A$ влияет на эффективность последующего применения $B$>>.

\paragraph{Вариант~2 (непосредственное соседство).} Воздействие фиксирует не только относительный порядок, но и непосредственное соседство:
\begin{equation*}
    T_{AB}^{(2)} = \begin{cases}
        1, & \text{если в } \pi \text{ присутствует подпоследовательность } [A, B]; \\
        0, & \text{в противном случае}.
    \end{cases}
\end{equation*}

Вариант~2 улавливает \emph{локальные синергии}: эффекты, возникающие только при непосредственном соседстве проходов. Доля прогонов с соседством $[A, B]$ существенно ниже (около 31~наблюдение на пару на бенчмарк, $\approx 16\,000$~наблюдений на пару суммарно), что снижает статистическую мощность по сравнению с вариантом~1, но даёт принципиально иной класс выявляемых зависимостей.

\paragraph{Зачем обе формулировки.} Каузальный вопрос <<влияет ли порядок>> не имеет единственного канонического воздействия: оба варианта описывают разные классы зависимостей, и совместное их использование позволяет различить глобальные и локальные эффекты. Конкретное сравнение результатов двух формулировок приведено в~\textbf{Главе~6}.

\subsection{Этапы каузального анализа}

Описанная ниже последовательность этапов применяется независимо к каждой из двух формулировок воздействия (1 и~2) и состоит из трёх обязательных шагов; четвёртый шаг (PC-проверка на медиаторы) рассмотрен как опциональное расширение и в настоящей работе не выполняется.

\subsubsection{Шаг 1. Вычисление ATE для всех пар проходов}

Для каждой из 3916~пар $(A, B)$:

\begin{enumerate}
    \item Из общего множества прогонов отбираются прогоны, в которых присутствуют оба прохода пары (см.~формулировку воздействия). Пары, для которых суммарный размер выборки воздействия и контроля составляет менее 20~наблюдений, исключаются из анализа как недостаточные для статистического вывода.
    \item Вычисляется оценка ATE: $\widehat{\text{ATE}}(A \to B) = \bar{Y}_1 - \bar{Y}_0$.
    \item Проверяется статистическая значимость разности средних с помощью теста Уэлча (раздел~3.4.1): вычисляется $t$-статистика и $p$-значение.
\end{enumerate}

Результат шага --- таблица из 3916~строк с полями: пара $(A, B)$, оценка ATE, $t$-статистика, $p$-значение, размер выборки воздействия~$n_1$, размер выборки контроля~$n_0$.

\subsubsection{Шаг 2. Коррекция множественного тестирования}

К полученному массиву из 3916 $p$-значений применяется процедура Бенджамини--Хохберга (раздел~3.4.2) с уровнем FDR~$\alpha = 0{,}05$. Результат --- скорректированные $p$-значения~$p_{\text{adj}}$ и индикатор значимости.

Бонферрони-коррекция не используется: при $K = 3916$ и $\alpha = 0{,}05$ она требует $p < 1{,}3 \cdot 10^{-5}$, что отсекает большинство реальных эффектов малой и средней величины.

\subsubsection{Шаг 3. Построение каузального графа}

По таблице ATE строится направленный граф зависимостей $G = (V, E)$, где $V$~--- множество из 89~проходов, а ребро $(A \to B) \in E$ включается, если для пары $(A, B)$ выполняются \emph{оба} условия:

\begin{enumerate}
    \item \textbf{Статистическая значимость}: $p_{\text{adj}} < 0{,}05$.
    \item \textbf{Практическая значимость}: $|\widehat{\text{ATE}}(A \to B)| \geqslant \delta$, где~$\delta = 0{,}005$.
\end{enumerate}

Направление ребра определяется знаком ATE: при $\widehat{\text{ATE}}(A \to B) > 0$ ребро направлено от $A$ к $B$ (порядок $A \to B$ лучше), при отрицательном знаке --- наоборот.

\paragraph{Обоснование порога $\delta = 0{,}005$.} Без порога~$\delta$ большое количество пар с микроскопическими, но статистически значимыми эффектами включается в граф. Эмпирический анализ показывает, что в построенной таблице ATE медианная величина $|d|$~Коэна для значимых пар составляет около $0{,}017$ --- в~10~раз меньше порога <<малого эффекта>> Коэна. Порог $\delta = 0{,}005$ (улучшение на 0{,}5~процентных пункта от baseline) был выбран после визуального анализа гистограммы $|\widehat{\text{ATE}}|$ значимых пар и соответствует <<минимально интересной>> практической величине эффекта для рассматриваемой метрики.

\subsection{Уровни агрегации}

Каузальный граф строится на трёх уровнях агрегации, отвечающих на разные исследовательские вопросы.

\subsubsection{Глобальный уровень}

ATE вычисляется по всем 517~бенчмаркам совокупно. Каждое наблюдение в выборке --- один прогон одного бенчмарка; оценка ATE отражает усреднённый каузальный эффект порядка пары проходов по всему ансамблю программ. Глобальный граф используется как опорное представление структуры зависимостей.

\subsubsection{Уровень доменных классов (per-class)}

Бенчмарки разбиваются на 8~доменных классов на основе ручной классификации по вычислительному характеру (раздел <<Группировки программ>> ниже): \texttt{crypto}, \texttt{numeric}, \texttt{multimedia}, \texttt{loop\_bench}, \texttt{pointer\_data}, \texttt{string\_parse}, \texttt{systems}, \texttt{general}. Для каждого класса ATE вычисляется только по принадлежащим ему бенчмаркам, и строится отдельный каузальный граф. Per-class графы отражают предметную специфику: например, эффект векторизации в классе \texttt{loop\_bench} (тесты циклов) ожидается более выраженным, чем в \texttt{systems} (системные приложения с преобладанием управления).

\subsubsection{Уровень структурных кластеров (per-cluster)}

Бенчмарки разбиваются на 6~кластеров по структуре промежуточного представления (см. раздел <<Группировки программ>>). Per-cluster графы отражают \emph{объективную} группировку по характеристикам IR --- в отличие от ручной классификации, не зависящую от знаний экспертизы. Сравнение per-class и per-cluster графов позволяет проверить, насколько домен программы определяет её каузальную структуру оптимизаций.

\medskip

Суммарно для каждой формулировки воздействия строится $1 + 8 + 6 = 15$~каузальных графов; для двух формулировок --- 30~графов.

\subsection{Группировки программ}

\subsubsection{Ручная классификация (per-class)}

Доменные классы введены на основе характера выполняемой программой задачи. Для бенчмарков из набора lac-dcc (231~программа) использована оригинальная классификация наборов: \texttt{cBench security\_*} отнесены к \texttt{crypto}, \texttt{PolyBench} --- к \texttt{numeric}, \texttt{TSVC} --- к \texttt{loop\_bench} и т.~д. Для функций из набора AnghaBench (286~функций) классификация выполнена индивидуально по имени файла и функции (имя содержит \texttt{extr\_<file>\_c\_<function>}), а не по проекту-источнику; такой подход даёт более точное отражение характера конкретной функции, в отличие от классификации по доменному характеру всего проекта. Финальное распределение приведено в~\textbf{Главе~5}.

\subsubsection{Автоматическая кластеризация (per-cluster)}

Структурная кластеризация выполняется в три шага.

\paragraph{Извлечение признаков IR.} Для каждого из 517~бенчмарков из \texttt{.bc} файла без применения оптимизаций (\texttt{baseline\_raw}) вычисляются 92~признака. Из них для кластеризации отбираются 21~\emph{нормированный} признак, не зависящий от размера программы:

\begin{itemize}
    \item Доли категорий инструкций (12~признаков): арифметические, логические, сравнения, обращения к памяти, переходы, вызовы и др.
    \item Доли ключевых опкодов (5): \texttt{load}, \texttt{store}, \texttt{alloca}, \texttt{getelementptr}, \texttt{br}.
    \item Структурные признаки (3): средний размер базового блока, средний размер функции, среднее число базовых блоков на функцию.
    \item Соотношения (1): \texttt{load\_store\_ratio}, \texttt{int\_float\_ratio}.
\end{itemize}

Абсолютные размеры (общее число инструкций, число функций) исключены, поскольку кластеризация по ним давала бы группировку <<большие>>/<<маленькие>> программы вместо группировки по структурным свойствам.

\paragraph{Снижение размерности и кластеризация.} Признаки стандартизованы (\texttt{StandardScaler}), затем применён метод главных компонент: оставлено 13~компонент, объясняющих 95{,}4~\% дисперсии. На пространстве компонент выполнена кластеризация методом $k$-средних.

\paragraph{Выбор числа кластеров.} Силуэтный коэффициент монотонно возрастает с увеличением~$k$ от $0{,}1455$ при $k=3$ до $0{,}1899$ при $k=15$, что свидетельствует об отсутствии чётко выраженной кластерной структуры (данные образуют непрерывное облако). Дополнительная проверка алгоритмом HDBSCAN подтверждает непрерывность: при минимальном размере кластера~25 выделяются лишь 3~кластера, и 75~\% программ помечаются как шум.

В этих условиях выбор~$k$ определяется содержательными соображениями. Зафиксировано~$k = 6$ как баланс между разрешением группировки и минимальным размером кластера: при~$k = 6$ минимальный кластер содержит 22~программы, что соответствует приблизительно 220\,000~наблюдений на пару (достаточно для статистической мощности per-cluster анализа). При~$k = 5$ один из кластеров содержал бы лишь 7~программ, при~$k = 7$~--- лишь 4.

\subsection{Условный средний эффект воздействия}

Для оценки гетерогенности каузальных эффектов по типам программ вычисляется условный средний эффект воздействия (раздел~3.2.4):
\begin{equation*}
    \text{CATE}(A \to B \mid g) = \mathbb{E}[Y \mid T_{AB} = 1, G = g] - \mathbb{E}[Y \mid T_{AB} = 0, G = g],
\end{equation*}
где $G$~--- принадлежность бенчмарка к группе (доменному классу или структурному кластеру).

CATE вычисляется только для рёбер \emph{глобального} каузального графа: для каждого из 608~значимых глобальных рёбер (см.~\textbf{Главу~6}) оценивается CATE по 8~доменным классам, 6~структурным кластерам и 517~отдельным бенчмаркам. На уровне отдельного бенчмарка для оценки CATE требуется не менее 10~наблюдений в подгруппе; подгруппы меньшего размера исключаются.

Результаты CATE используются для двух целей:
\begin{itemize}
    \item \textbf{Идентификация гетерогенности направления.} Если знак CATE на какой-либо подгруппе противоположен знаку глобального ATE, значит, ребро <<меняет направление>> на этой подгруппе. Доля таких рёбер количественно характеризует ограниченность глобального графа.
    \item \textbf{Выделение универсальных и специализированных рёбер.} Рёбра с однородным CATE по всем группам можно считать универсальными правилами; рёбра с существенной дисперсией CATE между группами требуют специализированных графов на уровне группы.
\end{itemize}

\subsection{Корреляционный граф для сравнения}

Для количественного сравнения каузального и корреляционного подходов на одних и тех же данных строится корреляционный граф. Корреляционный анализ применяется только к формулировке воздействия~2 (соседство), которая соответствует постановке корреляционного анализа через биграммы --- пары соседних проходов в последовательности.

\subsubsection{Определение корреляционного графа}

Для построения корреляционного графа исход бинаризуется: \emph{выигрышным} называется прогон, в котором применённая последовательность проходов даёт меньшее число IR-инструкций, чем стандартный пайплайн \texttt{-Oz}:
\begin{equation*}
    W = \mathbb{I}[M_\pi < M_{\text{-Oz}}],
\end{equation*}
где $M_{\text{-Oz}}$ --- число IR-инструкций после применения \texttt{-Oz} к тому же бенчмарку.

Для каждой упорядоченной пары проходов $(A, B)$ (всего $89 \cdot 88 = 7832$~пары) проверяется гипотеза о статистической независимости двух бинарных переменных: $T_{AB}^{(2)}$ (наличие подпоследовательности $[A, B]$ в~$\pi$) и $W$ (выигрышность прогона). Строится таблица сопряжённости 2$\times$2:
\begin{equation*}
    \begin{array}{c|cc}
        & W = 1 & W = 0 \\
        \hline
        T_{AB}^{(2)} = 1 & a & b \\
        T_{AB}^{(2)} = 0 & c & d
    \end{array}
\end{equation*}
и применяется критерий Пирсона~$\chi^2$ для проверки независимости:
\begin{equation*}
    \chi^2 = \sum_{ij} \frac{(O_{ij} - E_{ij})^2}{E_{ij}},
\end{equation*}
где $O_{ij}$ --- наблюдаемые частоты в ячейке таблицы, $E_{ij}$ --- ожидаемые частоты при выполнении гипотезы независимости.

Биграммы, встречающиеся менее чем в 20~прогонах, исключаются из анализа как недостаточные для построения таблицы сопряжённости. $p$-значения, полученные критерием~$\chi^2$ для оставшихся пар, корректируются процедурой Бенджамини--Хохберга на уровне FDR~$0{,}05$. Ребро $(A \to B)$ включается в корреляционный граф, если зависимость~$T_{AB}^{(2)}$ и~$W$ статистически значима и доля выигрышных прогонов с биграммой~$[A, B]$ превышает долю выигрышных среди прогонов без неё.

\subsubsection{Обоснование выбора метода}

Описанное определение корреляционного графа использует распространённую в работах по автотюнингу компиляторов схему: <<выигрышные>> последовательности (превосходящие фиксированный эталон \texttt{-Oz}) сопоставляются по частоте встречаемости отдельных переходов между проходами. Бинаризация исхода и анализ биграмм применяются в задачах автотюнинга начиная с работы Milepost GCC~\cite{fursin2011milepost}. Сравнение с подобным эталонным корреляционным анализом позволяет продемонстрировать выигрыш формального каузального вывода над подходом, фактически применяющимся на практике для построения корреляционных графов между проходами.

\subsubsection{Сравнение каузального и корреляционного графов}

Корреляционный граф строится на тех же уровнях агрегации, что и каузальный: глобальный, per-class и per-cluster. Сравнение проводится по следующим показателям:

\begin{itemize}
    \item \textbf{Число рёбер.} Превышение числа рёбер в корреляционном графе над числом рёбер в каузальном графе количественно характеризует масштаб ложных корреляций, индуцированных косвенными зависимостями.
    \item \textbf{Доля ложных корреляционных рёбер.} Доля рёбер корреляционного графа, отсутствующих в каузальном графе, интерпретируется как нижняя оценка доли корреляций без каузального содержания.
\end{itemize}

Полученные количественные оценки расхождения представлены в~\textbf{Главе~6} и используются в работе как центральный аргумент в пользу необходимости каузального анализа.

\newpage
